\section{Results}
\label{sec:results}

\subsection{Locked minority-class utility}

All 36 locked method comparisons completed (4 tasks $\times$ 3 generator seeds $\times$ 3 classifiers), with matched CTGAN pools, class quotas, and threshold policies. Table~\ref{tab:locked-primary} gives the primary results. The prespecified operating-regime rule passed: both Vigitel tasks had positive minority-F1 and Macro-F1 differences, at least 6 of 9 paired-model wins, and Holm-adjusted primary $p<.05$.

\begin{table}[htbp]
\centering
\caption{Internal locked TSTR results. Differences are \sepa{} minus matched Standard CTGAN. Confidence intervals are feature-group bootstrap intervals; $p$ values are Holm-adjusted across the four primary minority-F1 comparisons.}
\label{tab:locked-primary}
\small
\begin{tabularx}{\textwidth}{Xrrrrrr}
\toprule
Task & Standard F1 & \sepa{} F1 & Difference (95\% CI) & Holm $p$ & Wins/9 & $\Delta$ECE\\
\midrule
COVID-19 mortality & .341 & .347 & .006 ($-.121$, .131) & .925 & 5 & .057\\
Kidney mortality & .062 & .118 & .056 (.013, .107) & .057 & 8 & .038\\
Vigitel obesity & .199 & .311 & .112 (.074, .150) & $<.001$ & 9 & .099\\
Vigitel smoking & .134 & .276 & .143 (.096, .190) & $<.001$ & 9 & .085\\
\bottomrule
\end{tabularx}
\end{table}


For Vigitel obesity, minority-F1 increased from 0.199 to 0.311 (difference 0.112, 95\% CI 0.074--0.150), and for current smoking it increased from 0.134 to 0.276 (difference 0.143, 95\% CI 0.096--0.190). Both effects occurred in all nine seed--classifier pairs. The corresponding Macro-F1 differences were 0.067 and 0.090, and AUPRC differences were 0.082 and 0.130.

COVID-19 minority-F1 changed by 0.006 with a wide interval spanning material benefit and harm. Its Macro-F1 increased by 0.089, but only five of nine minority-F1 pairs favored \sepa{}; the primary classification claim was therefore not supported. Kidney minority-F1 increased by 0.056 and favored \sepa{} in eight of nine pairs, but the four-test Holm adjustment gave $p=.057$. Kidney Macro-F1 changed by only 0.009. These two outcomes are boundary conditions rather than required wins.

\subsection{Decision thresholds and calibration}

The learned thresholds differed from 0.5 in every task. Mean Standard/\sepa{} thresholds were 0.221/0.307 for COVID-19, 0.085/0.126 for kidney mortality, 0.187/0.255 for obesity, and 0.153/0.281 for smoking. These values show that synthetic-to-real classification involved a substantial probability-scale shift even though threshold choice never used real labels.

Classification gains did not imply calibrated risk estimates. Brier score worsened by 0.031 for obesity and 0.036 for smoking; ECE worsened by 0.099 and 0.085. ECE also increased for COVID-19 and kidney mortality. The confirmed Vigitel result is therefore an improvement in minority detection at a training-derived operating point, not evidence of probability calibration.

\subsection{Selection mechanism}

Selected rows differed substantially from quota-matched Standard CTGAN rows. Before diversity repair, minority-selection overlap was close to random expectation in all four tasks. Progressive replacement changed held-out utility while the generator and pool remained fixed (Table~\ref{tab:mechanism}). For Vigitel obesity and smoking, replacing Standard rows with curated rows increased TSTR minority-F1 by 0.132 and 0.174, respectively. COVID-19 also showed a positive development-stage replacement effect, but this did not reproduce as a locked primary effect. Kidney showed partial selection benefit under a strong generator and diversity constraint.

\begin{table}[htbp]
\centering
\caption{Saved-pool mechanism evidence. Pool coverage is the fraction of real positive profiles reached by the shared candidate pool. Replacement effects are changes after progressively substituting Standard-selected with \sepa{}-selected rows; they are mechanism estimates on development partitions, not locked confirmation.}
\label{tab:mechanism}
\small
\begin{tabularx}{\textwidth}{Xrrrr}
\toprule
Task & Pool coverage & Selection Jaccard & $\Delta$Macro-F1 & $\Delta$minority-F1\\
\midrule
COVID-19 mortality & .971 & .078 & .080 & .077\\
Kidney mortality & .969 & .402 & $-.010$ & .082\\
Vigitel obesity & .894 & .084 & .084 & .132\\
Vigitel smoking & .988 & .019 & .107 & .174\\
\bottomrule
\end{tabularx}
\end{table}


The kidney result clarifies the role of the diversity repair. Pure Pareto selection shared 1.03\% of minority rows with Standard CTGAN but retained only 57.6\% of Standard diversity. The frozen repair consequently restored Standard-selected rows in every kidney block, replacing an average of 55\% of the final minority subset and raising final overlap to 55.5\%. The final kidney method was therefore a heavily repaired hybrid, which limited the contrast between training sets.

\subsection{Coverage and distributional trade-offs}

Across the five saved diagnostic folds, curation increased local $k$-nearest-neighbor precision from 0.611 to 0.988 for COVID-19, 0.315 to 0.552 for kidney mortality, 0.437 to 0.870 for obesity, and 0.599 to 0.803 for smoking. At the same time, $k$-nearest-neighbor recall decreased from 0.943 to 0.804, 0.898 to 0.789, 0.839 to 0.729, and 0.874 to 0.727, respectively. Rare-mode coverage decreased in all four tasks.

This pattern is consistent with prototype concentration: the selected cases are more locally plausible, but they span less of the positive-class support. The Vigitel gains occurred despite worse coverage, so broader support did not mediate their classification improvement. In kidney mortality, coverage of high erythropoietin-count profiles fell from 0.770 under Standard CTGAN to 0.558 under \sepa{}, although the shared pool reached 0.816. Some clinically relevant trajectories were therefore available at generation but lost during curation.

\subsection{Generator-budget sensitivity}

The epoch study showed that curation advantage and generator maturity are not interchangeable. Thirty epochs were sufficient for the tested COVID-19 endpoints and did not explain its instability. Longer kidney training did not rescue augmentation and produced substantial duplication and exact-training-row matches at 300--1,024 epochs. For Vigitel obesity, Standard CTGAN improved strongly by 300 epochs, narrowing the curation advantage. For smoking TSTR, Standard Macro-F1 increased from 0.486 at 30 epochs to 0.605 at 1,024 epochs and matched \sepa{} at 1,024 (difference $-0.001$). In smoking augmentation, a smaller \sepa{} advantage persisted.

Thus, part of the low-budget effect is efficient recovery of structure already latent in an imperfect candidate pool. It should not be interpreted as unconditional superiority over a fully trained generator. Conversely, additional training did not solve the COVID-19 information constraint or the kidney sparse-trajectory problem.
